\section{Từ lời giải thích cục bộ đến một kết luận toàn cục}
\label{sec:ch04-tu-cuc-bo-den-toan-cuc}

\index{suy luận thống kê}%
Vì TreeSHAP chỉ làm phép tính cục bộ nhanh hơn, một kết luận toàn cục vẫn phải
được xây từ nhiều đầu vào. Với đặc trưng $i$ cố định, $\phi_i(x)$ là một hàm
theo $x$: nó cho phần đóng góp của $i$ tại từng điểm. Trong thực hành, người
ta thường lấy trung bình độ lớn hoặc bình phương của hàm này trên một
tập dữ liệu để xếp hạng đặc trưng. Phép tổng hợp ấy chuyển một lời giải thích
cục bộ thành một phát biểu toàn cục, nên nó cần một câu hỏi mới: độ bất định
của con số tổng hợp là bao nhiêu?

Whitehouse, Sawarni và Syrgkanis gọi hàm $\phi_i(\cdot)$ ấy là đường cong
SHAP, và xét các lũy thừa bậc $p$ của nó; trung bình trị tuyệt đối ứng với
$p=1$ và trung bình bình phương ứng với $p=2$. Họ coi đường cong SHAP là một
hàm phụ phải ước lượng từ dữ liệu, nên thay một mô hình đã học vào rồi lấy
trung bình thì chưa đủ để có suy luận tiệm cận chuẩn. Bài báo xây dựng các ước
lượng đã khử chệch để lập khoảng tin cậy cho các đại lượng
đó~\autocite{p19shapinference}.

\index{lời giải thích cục bộ}%
Kết quả này không biến một giá trị SHAP cục bộ thành bằng chứng về cơ chế của
mô hình. Nó làm một việc khác: nếu ta đã chọn hàm giá trị và dữ liệu đánh giá,
thì phát biểu toàn cục dựa trên các giá trị SHAP phải mang theo bất định lấy
mẫu và bất định do ước lượng. Đây là nền tảng thống kê mà một biểu đồ xếp hạng
đặc trưng thường chưa thể hiện.
